% Explanation certificate for heart-s4-c1 (genome c85804fa-08bd-4207-ac69-a31a8567473c).
% Generated by explaneat export (schema 1.4.0) from the JSON record alone;
% content sha256 5bb264e6720c1a25b94714c430f44b3b1f28f41592eef373c0a2d422a381b5a0.
% Drafter recorded in the JSON: ai_agent (claude). Do not edit by hand: re-render with
%   uv run python -m explaneat export render --from heart-s4-c1.json
% TikZ modules are input from \explaneatcertdir{<label>} (empty by default: same
% directory; index.tex sets it to the genome's subdirectory).
\providecommand{\explaneatcertdir}[1]{}
\begin{drafted}[appendix: heart-s4-c1 explanation certificate]
\section{Explanation certificate: heart-s4-c1}\label{app:heart-s4-c1:sec}
This section is the explanation certificate of the heart-s4-c1 network. It is rendered from the canonical JSON record \texttt{heart-s4-c1.json} (schema 1.4.0, SHA-256 \texttt{5bb264e6720c1a25}\ldots), which also holds every evidence payload, the full graphs with weights, and the explanation record. \Cref{app:heart-s4-c1:summary} summarises the network and its coverage.

\begin{table}[htbp]
\centering\small
\begin{tabularx}{\linewidth}{@{}lX@{}}
\toprule
Dataset & Heart Disease (binary, prepared) \\
Experiment & 10seed-heart-s4, seed 4 \\
Genome & \texttt{c85804fa-08bd-4207-ac69-a31a8567473c} (NEAT key 9679, generation 39) \\
Explanation & \texttt{10565a81-f55b-47f4-b51e-21641ec08556} (36 operations) \\
Train AUC / accuracy & 0.9423 / 0.8884 (n = 242) \\
Test AUC / accuracy & 0.8858 / 0.8197 (n = 61) \\
Output activation & sigmoid \\
Phenotype & 31 nodes (26 inputs, 4 hidden, 1 output), 29 enabled edges \\
Explained model state & 36 nodes, 29 enabled edges \\
Refactoring operations & split\_\allowbreak{}node 5 \\
Function preservation & max $|\Delta\,\mathrm{output}|$ = $3.33\times 10^{-16}$ on the training split \\
Annotations & 5 (4 leaf, 1 composition) \\
Structural coverage & 1 (35 of 35 candidate nodes; 6 Null) \\
Compositional coverage & 1 \\
Uncovered nodes / edges & 0 / 0 \\
Missing compositions & 0 \\
Full coverage & yes \\
Drafted by & AI-assisted draft \\
Verified by & unverified \\
Record hash (SHA-256) & \texttt{5bb264e6720c1a25}\ldots \\
\bottomrule
\end{tabularx}
\caption{Explanation certificate summary: heart-s4-c1.}
\label{app:heart-s4-c1:summary}
\end{table}

\paragraph{Drafter's summary.} Essentially logistic regression: 19 direct edges (logit SD 4.55; train AUC 0.933 alone; ca +2.25, slope:2 +1.12, sex:1 +0.90) plus four additive hidden readouts (logit SD 0.86) to 0.942. The one meaningful nonlinearity is an AND gate 2065 = 3.64*[thal:6 AND fbs=0] read out -1.08, lowering the logit by 3.9 for the 10 fixed-defect patients with normal fasting sugar (core mean predicted 0.83, disease rate 50\%). The cp:4 and slope:3 hidden indicators are exactly affine on their binary inputs and a floored thalach relay is near-inert. Same gate and linear core as the 10-seed heart-s4 record (key 9650, gen 38); here cp:4 and slope:3 indicators replace that genome's cp:2 indicator.

\paragraph{What resisted explanation.} Nothing resisted coverage. The one real interaction (fixed thallium defect AND normal fasting sugar, gate 2065) rests on 10 training patients and may be partly noise-fitted. Hidden nodes 1754 (cp:4) and 3504 (slope:3) are exactly affine on their binary inputs, so they are re-parameterisations of linear terms, not mechanisms. Counter-directional signs reported, not explained: fbs:1 -0.58 in the linear core, and the gate making a fixed defect protective for normal-sugar patients.

\subsection{Plain-language summary}\label{app:heart-s4-c1:plain}
\noindent{\small\textit{AI-assisted draft for review (unverified). Each sentence cites the claims (C01, C02, \ldots) that back it; every cited claim was computed on the model and the training split and passed its automatic checks. Magnitudes are changes in predicted risk for a typical patient; confidence words follow the test-split check and n. The text describes what the model does, not the world.}}\par\smallskip
\paragraph{Summary (drafted).} For a typical patient (a 56-year-old man with no chest pain symptoms and median test results) the model predicts a risk of heart disease of about 24\%; it adds up most of its inputs independently, and four small parts add to that sum: one for a fixed defect on the thallium scan with normal fasting blood sugar, one for chest pain type, one for a downsloping ST segment and one for maximum heart rate.\,{\scriptsize\textcolor{NavyBlue}{[C31]}} These add up independently, and these lower the risk: being a woman takes a typical patient's predicted risk from about 24\% to about 1\%; through this sum alone, compared with no chest pain symptoms, typical angina lowers it to about 2\%, non-anginal chest pain to about 3\% and atypical angina to about 14\%, fasting blood sugar above 120 mg/dL lowers it to about 7\%, and a higher maximum heart rate lowers it steadily, from about 27\% at 118 to about 21\% at 175 beats per minute; chest pain type, fasting blood sugar and maximum heart rate also act through the parts below.\,{\scriptsize\textcolor{NavyBlue}{[C12, C13, C15, C22, C17, C24]}} Each extra vessel coloured by fluoroscopy adds about the same, from about 24\% at 0 to about 98\% at 2 vessels, flattening out near the top because the predicted risk cannot go much higher; through this sum alone a flat rather than upsloping ST segment raises it to about 82\% and any thallium result other than a normal scan to about 48\%, and chest pain brought on by exercise raises it to about 36\% (the ST segment and the thallium scan also act through the parts below).\,{\scriptsize\textcolor{NavyBlue}{[C01, C02, C06, C21]}} Higher cholesterol (from about 14\% at 193 mg/dL to about 44\% at 308 mg/dL), more ST depression on exercise (from about 18\% at 0 mm to about 42\% at 2.8 mm), higher resting blood pressure (from about 14\% at 110 mm Hg to about 41\% at 155 mm Hg) and older age (from about 21\% at 42 to about 25\% at 66 years) raise it steadily.\,{\scriptsize\textcolor{NavyBlue}{[C04, C08, C05, C25]}} Through this sum alone a downsloping ST segment raises it to about 39\% (the model has little data on this, n = 17), and an ST-T wave abnormality on the resting ECG raises it to about 38\% (the model has little data on this, n = 2).\,{\scriptsize\textcolor{NavyBlue}{[C19, C20]}} This part reads fasting blood sugar and the thallium scan and acts on 10 of 242 training patients, nearly all of them with a fixed defect on the thallium scan, 5 of whom had heart disease; for the typical patient it acts on, the predicted risk is about 94\% without it and about 24\% with it (the model has little data on this, n = 10).\,{\scriptsize\textcolor{NavyBlue}{[C35]}} Through it, a fixed defect on the thallium scan lowers a typical patient's predicted risk from about 24\% to about 2\% if fasting blood sugar is 120 mg/dL or less, but the other way if it is above 120 mg/dL, from about 7\% to about 18\% (the model has little data on this, n = 15).\,{\scriptsize\textcolor{NavyBlue}{[C34]}} For the typical patient it acts on, fasting blood sugar above 120 mg/dL raises the risk through it from about 24\% to about 94\%.\,{\scriptsize\textcolor{NavyBlue}{[C32]}} This part reads only chest pain type and acts on the 134 of 242 training patients who have some chest pain symptom, 30 of whom had heart disease; for the typical patient it acts on, the predicted risk is about 2\% without it and about 1\% with it.\,{\scriptsize\textcolor{NavyBlue}{[C38]}} Through it, any chest pain type other than no chest pain symptoms lowers a typical patient's predicted risk from about 24\% to about 12\%.\,{\scriptsize\textcolor{NavyBlue}{[C37]}} This part reads only the slope of the ST segment and acts on the 17 of 242 training patients with a downsloping ST segment, 8 of whom had heart disease; for the typical patient it acts on, the predicted risk is about 66\% without it and about 73\% with it (the model has little data on this, n = 17).\,{\scriptsize\textcolor{NavyBlue}{[C41]}} Through it, a downsloping ST segment raises a typical patient's predicted risk from about 24\% to about 29\% (the model has little data on this, n = 17).\,{\scriptsize\textcolor{NavyBlue}{[C40]}} This part reads only maximum heart rate and acts on 190 of 242 training patients, 71 of whom had heart disease; for the typical patient it acts on, the predicted risk is about 21\% without it and about 20\% with it.\,{\scriptsize\textcolor{NavyBlue}{[C45]}} Through it, a higher maximum heart rate changes nothing until about 131 beats per minute, then lowers the predicted risk steadily, for a typical patient from about 24\% at 118 to about 23\% at 175 beats per minute.\,{\scriptsize\textcolor{NavyBlue}{[C43, C44]}}
\paragraph{The whole model.} The model adds up its inputs into one score: the vessels coloured (each extra vessel adding about the same), the ST segment, sex and chest pain type each move the predicted risk on their own, and four small parts add to that sum: one for a fixed thallium defect with normal fasting blood sugar, one for chest pain type, one for a downsloping ST segment and one for maximum heart rate.\,{\scriptsize\textcolor{NavyBlue}{[C31, C01, C02, C12, C10]}} \textit{Evidence.} It is made up of the inputs that feed the score directly and four parts, the fixed thallium defect part, the chest pain offset, the downsloping ST offset and the maximum heart rate relay (each explained below); the score is their sum, which was checked on all 242 training patients.\,{\scriptsize\textcolor{NavyBlue}{[C31]}} Changing one input at a time across the whole model, compared with no chest pain symptoms, typical angina lowers a typical patient's predicted risk from about 24\% to about 1\%, non-anginal chest pain to about 1\% and atypical angina to about 7\%, against about 2\% for typical angina through the sum alone.\,{\scriptsize\textcolor{NavyBlue}{[C10, C11, C16, C13]}} Across the whole model, fasting blood sugar above 120 mg/dL lowers it from about 24\% to about 7\%, the same as through the sum alone.\,{\scriptsize\textcolor{NavyBlue}{[C18, C17]}} Across the whole model, a higher maximum heart rate lowers it slowly at first, then faster, from about 28\% at 118 to about 21\% at 175 beats per minute, against about 27\% to about 21\% through the sum alone, where it falls steadily.\,{\scriptsize\textcolor{NavyBlue}{[C23, C24]}} Across the whole model, a downsloping rather than upsloping ST segment raises it to about 47\% (the model has little data on this, n = 17), and a reversible defect on the thallium scan or a thallium result coded 4 raises it to about 48\%.\,{\scriptsize\textcolor{NavyBlue}{[C09, C07]}}
\paragraph{Fixed thallium defect with normal sugar.} This part lowers the predicted risk for patients with a fixed defect on the thallium scan when their fasting blood sugar is 120 mg/dL or less (the model has little data on this, n = 15).\,{\scriptsize\textcolor{NavyBlue}{[C34, C35]}} \textit{Evidence.} It acts on 10 of 242 training patients, nearly all of them with a fixed defect on the thallium scan, 5 of whom had heart disease; for the typical patient it acts on, the predicted risk is about 94\% without it and about 24\% with it (the model has little data on this, n = 10).\,{\scriptsize\textcolor{NavyBlue}{[C35]}} Through it, a fixed defect on the thallium scan lowers a typical patient's predicted risk from about 24\% to below 1\% (the model has little data on this, n = 15).\,{\scriptsize\textcolor{NavyBlue}{[C33]}} Through it, a fixed defect lowers the risk from about 24\% to about 2\% if fasting blood sugar is 120 mg/dL or less, but raises it from about 7\% to about 18\% if it is above 120 mg/dL (the model has little data on this, n = 15).\,{\scriptsize\textcolor{NavyBlue}{[C34]}} For the typical patient it acts on, fasting blood sugar above 120 mg/dL raises the risk through it from about 24\% to about 94\%.\,{\scriptsize\textcolor{NavyBlue}{[C32]}} Its formula reproduces the part exactly on all 242 training patients.\,{\scriptsize\textcolor{NavyBlue}{[C36]}}
\paragraph{Chest pain offset.} This part lowers the predicted risk for patients with any chest pain symptom, compared with patients who have none.\,{\scriptsize\textcolor{NavyBlue}{[C37]}} \textit{Evidence.} It acts on the 134 of 242 training patients who have some chest pain symptom, 30 of whom had heart disease; for the typical patient it acts on, the predicted risk is about 2\% without it and about 1\% with it.\,{\scriptsize\textcolor{NavyBlue}{[C38]}} Through it, any chest pain type other than no chest pain symptoms takes a typical patient's predicted risk from about 24\% to about 12\%.\,{\scriptsize\textcolor{NavyBlue}{[C37]}} Its formula reproduces the part exactly on all 242 training patients.\,{\scriptsize\textcolor{NavyBlue}{[C39]}}
\paragraph{Downsloping ST offset.} This part raises the predicted risk a little for patients with a downsloping ST segment (the model has little data on this, n = 17).\,{\scriptsize\textcolor{NavyBlue}{[C40]}} \textit{Evidence.} It acts on the 17 of 242 training patients with a downsloping ST segment, 8 of whom had heart disease; for the typical patient it acts on, the predicted risk is about 66\% without it and about 73\% with it (the model has little data on this, n = 17).\,{\scriptsize\textcolor{NavyBlue}{[C41]}} Through it, a downsloping ST segment takes a typical patient's predicted risk from about 24\% to about 29\% (the model has little data on this, n = 17).\,{\scriptsize\textcolor{NavyBlue}{[C40]}} Its formula reproduces the part exactly on all 242 training patients.\,{\scriptsize\textcolor{NavyBlue}{[C42]}}
\paragraph{Maximum heart rate relay.} Through this part, a higher maximum heart rate changes nothing until about 131 beats per minute, then lowers the predicted risk steadily.\,{\scriptsize\textcolor{NavyBlue}{[C43, C44]}} \textit{Evidence.} It acts on 190 of 242 training patients, 71 of whom had heart disease; for the typical patient it acts on, the predicted risk is about 21\% without it and about 20\% with it.\,{\scriptsize\textcolor{NavyBlue}{[C45]}} Through it, a higher maximum heart rate changes nothing until about 131 beats per minute, then lowers the risk steadily, for a typical patient from about 24\% at 118 to about 23\% at 175 beats per minute.\,{\scriptsize\textcolor{NavyBlue}{[C43, C44]}} Its formula reproduces the part exactly on all 242 training patients.\,{\scriptsize\textcolor{NavyBlue}{[C47]}}
\paragraph{Worth checking against what is known.} Worth checking: going from 0 to 2 vessels coloured moves a typical patient's predicted risk from about 24\% to about 98\%.\,{\scriptsize\textcolor{NavyBlue}{[C01]}} Worth checking: patients with typical angina get the lowest predicted risk of all chest pain types, below patients with no chest pain symptoms.\,{\scriptsize\textcolor{NavyBlue}{[C10, C11, C16]}} Worth checking: the model treats fasting blood sugar above 120 mg/dL as lowering the risk of heart disease (for a typical patient from about 24\% to about 7\%).\,{\scriptsize\textcolor{NavyBlue}{[C18]}} Worth checking: through one part, the model treats a fixed defect on the thallium scan as lowering the risk for patients with normal fasting blood sugar (for a typical patient from about 24\% to about 2\%), from only 15 training patients.\,{\scriptsize\textcolor{NavyBlue}{[C34]}}
\paragraph{Claims.}
\par{\tiny\noindent
\begin{tabularx}{\linewidth}{@{}l>{\raggedright\arraybackslash}X>{\raggedright\arraybackslash}p{0.11\linewidth}>{\raggedright\arraybackslash}p{0.16\linewidth}>{\raggedright\arraybackslash}p{0.05\linewidth}>{\raggedright\arraybackslash}p{0.15\linewidth}l>{\raggedright\arraybackslash}p{0.08\linewidth}@{}}
\toprule
Claim & What is checked & Risk, typical patient & Estimate (log-odds; OR); range across patients & n & Test check & Verified & Words \\
\midrule
C01 & direction: number of vessels coloured by fluoroscopy: 0 vessels to 2 vessels (whole model) \newline\textit{shape: steady (equal steps); 100\% of patients same; verified; risk saturates near 100\%} & 24\% to 98\% & log-odds 5.04, OR 154; range across patients 5.04 to 5.04 & 242 & train; test unclear (AUC gain 0.044 [-0.038, 0.127], n=61) & \textcolor{OliveGreen}{yes} & on the training data only \\
C02 & direction: slope of the ST segment at peak exercise: an upsloping ST segment to a flat ST segment (direct contribution only, not through any part) \newline\textit{shape: categorical; highest to lowest: a flat ST segment, a downsloping ST segment, an upsloping ST segment} & 24\% to 82\% & log-odds 2.72, OR 15.1; range across patients 2.72 to 2.72 & 115, 110 & train; test disagrees (AUC gain -0.014 [-0.061, 0.025], n=61) & \textcolor{OliveGreen}{yes} & on the training data only \\
C03 & direction: slope of the ST segment at peak exercise: an upsloping ST segment to a flat ST segment (whole model) \newline\textit{shape: categorical; highest to lowest: a flat ST segment, a downsloping ST segment, an upsloping ST segment} & 24\% to 82\% & log-odds 2.72, OR 15.1; range across patients 2.72 to 2.72 & 115, 110 & train; test disagrees (AUC gain -0.013 [-0.058, 0.026], n=61) & \textcolor{OliveGreen}{yes} & on the training data only \\
C04 & direction: serum cholesterol: 193 mg/dL (5.0 mmol/L) to 308 mg/dL (8.0 mmol/L) (whole model) \newline\textit{shape: steady; 100\% of patients same; verified} & 14\% to 44\% & log-odds 1.60, OR 4.97; range across patients 1.60 to 1.60 & 242 & train; test disagrees (AUC gain -0.019 [-0.049, 0.002], n=61) & \textcolor{OliveGreen}{yes} & on the training data only \\
C05 & direction: resting blood pressure: 110 mm Hg to 155 mm Hg (whole model) \newline\textit{shape: steady; 100\% of patients same; verified} & 14\% to 41\% & log-odds 1.48, OR 4.37; range across patients 1.48 to 1.48 & 242 & train; test unclear (AUC gain 0.002 [-0.018, 0.023], n=61) & \textcolor{OliveGreen}{yes} & on the training data only \\
C06 & direction: thallium stress test result: a normal thallium scan to thallium stress test result other than a normal thallium scan (direct contribution only, not through any part) \newline\textit{shape: categorical; highest to lowest: thallium stress test result other than a normal thallium scan, a normal thallium scan} & 24\% to 48\% & log-odds 1.09, OR 2.98; range across patients 1.09 to 1.09 & 135, 107 & train; test agrees (AUC gain 0.026 [0.005, 0.055], n=61) & \textcolor{OliveGreen}{yes} & clearly \\
C07 & direction: thallium stress test result: a normal thallium scan to thallium stress test result other than a normal thallium scan, a fixed defect on the thallium scan (whole model) \newline\textit{shape: categorical; highest to lowest: thallium stress test result other than a normal thallium scan, a fixed defect on the thallium scan, a normal thallium scan, a fixed defect on the thallium scan} & 24\% to 48\% & log-odds 1.09, OR 2.98; range across patients 1.09 to 1.09 & 135, 92 & train; test unclear (AUC gain 0.005 [-0.028, 0.039], n=61) & \textcolor{OliveGreen}{yes} & on the training data only \\
C08 & direction: ST depression on exercise: 0.0 mm to 2.8 mm (whole model) \newline\textit{shape: steady; 100\% of patients same; verified} & 18\% to 42\% & log-odds 1.16, OR 3.19; range across patients 1.16 to 1.16 & 242 & train; test unclear (AUC gain 0.005 [-0.01, 0.024], n=61) & \textcolor{OliveGreen}{yes} & on the training data only \\
C09 & direction: slope of the ST segment at peak exercise: an upsloping ST segment to a downsloping ST segment (whole model) \newline\textit{shape: categorical; highest to lowest: a flat ST segment, a downsloping ST segment, an upsloping ST segment} & 24\% to 47\% & log-odds 1.05, OR 2.85; range across patients 1.05 to 1.05 & 115, 17 (small n) & train; test disagrees (AUC gain -0.013 [-0.058, 0.026], n=61) & \textcolor{OliveGreen}{yes} & the model has little data on this; on the training data only \\
C10 & direction: chest pain type: no chest pain (asymptomatic) to typical angina (whole model) \newline\textit{shape: categorical; highest to lowest: no chest pain (asymptomatic), atypical angina, non-anginal chest pain, typical angina} & 24\% to 1\% & log-odds -3.84, OR 0.02; range across patients -3.84 to -3.84 & 108, 22 & train; test agrees (AUC gain 0.057 [0.016, 0.105], n=61) & \textcolor{OliveGreen}{yes} & clearly \\
C11 & direction: chest pain type: no chest pain (asymptomatic) to non-anginal chest pain (whole model) \newline\textit{shape: categorical; highest to lowest: no chest pain (asymptomatic), atypical angina, non-anginal chest pain, typical angina} & 24\% to 1\% & log-odds -3.12, OR 0.04; range across patients -3.12 to -3.12 & 108, 70 & train; test agrees (AUC gain 0.057 [0.016, 0.105], n=61) & \textcolor{OliveGreen}{yes} & clearly \\
C12 & direction: sex: male to female (whole model) \newline\textit{shape: categorical; highest to lowest: male, female} & 24\% to 1\% & log-odds -3.07, OR 0.05; range across patients -3.07 to -3.07 & 159, 83 & train; test unclear (AUC gain 0.025 [-0.01, 0.069], n=61) & \textcolor{OliveGreen}{yes} & on the training data only \\
\bottomrule
\end{tabularx}}
\par{\tiny\noindent
\begin{tabularx}{\linewidth}{@{}l>{\raggedright\arraybackslash}X>{\raggedright\arraybackslash}p{0.11\linewidth}>{\raggedright\arraybackslash}p{0.16\linewidth}>{\raggedright\arraybackslash}p{0.05\linewidth}>{\raggedright\arraybackslash}p{0.15\linewidth}l>{\raggedright\arraybackslash}p{0.08\linewidth}@{}}
\toprule
Claim & What is checked & Risk, typical patient & Estimate (log-odds; OR); range across patients & n & Test check & Verified & Words \\
\midrule
C13 & direction: chest pain type: no chest pain (asymptomatic) to typical angina (direct contribution only, not through any part) \newline\textit{shape: categorical; highest to lowest: no chest pain (asymptomatic), atypical angina, non-anginal chest pain, typical angina} & 24\% to 2\% & log-odds -2.98, OR 0.05; range across patients -2.98 to -2.98 & 108, 22 & train; test unclear (AUC gain 0.029 [0.000, 0.065], n=61) & \textcolor{OliveGreen}{yes} & on the training data only \\
C14 & direction: thallium stress test result: a normal thallium scan to a fixed defect on the thallium scan (whole model) \newline\textit{shape: categorical; highest to lowest: thallium stress test result other than a normal thallium scan, a fixed defect on the thallium scan, a normal thallium scan, a fixed defect on the thallium scan} & 24\% to 2\% & log-odds -2.29, OR 0.10; range across patients -2.83 to 1.09 & 135, 15 (small n) & train; test unclear (AUC gain 0.005 [-0.028, 0.039], n=61) & \textcolor{BrickRed}{no} & the model has little data on this; on the training data only \\
C15 & direction: chest pain type: no chest pain (asymptomatic) to non-anginal chest pain (direct contribution only, not through any part) \newline\textit{shape: categorical; highest to lowest: no chest pain (asymptomatic), atypical angina, non-anginal chest pain, typical angina} & 24\% to 3\% & log-odds -2.27, OR 0.10; range across patients -2.27 to -2.27 & 108, 70 & train; test unclear (AUC gain 0.029 [0.000, 0.065], n=61) & \textcolor{OliveGreen}{yes} & on the training data only \\
C16 & direction: chest pain type: no chest pain (asymptomatic) to atypical angina (whole model) \newline\textit{shape: categorical; highest to lowest: no chest pain (asymptomatic), atypical angina, non-anginal chest pain, typical angina} & 24\% to 7\% & log-odds -1.47, OR 0.23; range across patients -1.47 to -1.47 & 108, 42 & train; test agrees (AUC gain 0.057 [0.016, 0.105], n=61) & \textcolor{OliveGreen}{yes} & clearly \\
C17 & direction: fasting blood sugar: fasting blood sugar of 120 mg/dL or less to fasting blood sugar above 120 mg/dL (direct contribution only, not through any part) \newline\textit{shape: categorical; highest to lowest: fasting blood sugar of 120 mg/dL or less, fasting blood sugar above 120 mg/dL} & 24\% to 7\% & log-odds -1.43, OR 0.24; range across patients -1.43 to -1.43 & 209, 33 & train; test disagrees (AUC gain -0.002 [-0.020, 0.015], n=61) & \textcolor{OliveGreen}{yes} & on the training data only \\
C18 & direction: fasting blood sugar: fasting blood sugar of 120 mg/dL or less to fasting blood sugar above 120 mg/dL (whole model) \newline\textit{shape: categorical; highest to lowest: fasting blood sugar of 120 mg/dL or less, fasting blood sugar above 120 mg/dL} & 24\% to 7\% & log-odds -1.18, OR 0.31; range across patients -1.43 to -1.43 & 209, 33 & train; test disagrees (AUC gain -0.005 [-0.025, 0.013], n=61) & \textcolor{OliveGreen}{yes} & on the training data only \\
C19 & direction: slope of the ST segment at peak exercise: an upsloping ST segment to a downsloping ST segment (direct contribution only, not through any part) \newline\textit{shape: categorical; highest to lowest: a flat ST segment, a downsloping ST segment, an upsloping ST segment} & 24\% to 39\% & log-odds 0.75, OR 2.11; range across patients 0.75 to 0.75 & 115, 17 (small n) & train; test disagrees (AUC gain -0.014 [-0.061, 0.025], n=61) & \textcolor{OliveGreen}{yes} & the model has little data on this; on the training data only \\
C20 & direction: resting ECG: resting ECG other than an ST-T wave abnormality on the resting ECG to an ST-T wave abnormality on the resting ECG (whole model) \newline\textit{shape: categorical; highest to lowest: an ST-T wave abnormality on the resting ECG, resting ECG other than an ST-T wave abnormality on the resting ECG} & 24\% to 38\% & log-odds 0.68, OR 1.98; range across patients 0.68 to 0.68 & 240, 2 (small n) & train; test unclear (AUC gain 0.000 [0.000, 0.000], n=61) & \textcolor{OliveGreen}{yes} & the model has little data on this; on the training data only \\
C21 & direction: exercise-induced angina: no chest pain on exercise to chest pain brought on by exercise (whole model) \newline\textit{shape: categorical; highest to lowest: chest pain brought on by exercise, no chest pain on exercise} & 24\% to 36\% & log-odds 0.60, OR 1.81; range across patients 0.60 to 0.60 & 160, 82 & train; test unclear (AUC gain 0.009 [-0.002, 0.026], n=61) & \textcolor{OliveGreen}{yes} & on the training data only \\
C22 & direction: chest pain type: no chest pain (asymptomatic) to atypical angina (direct contribution only, not through any part) \newline\textit{shape: categorical; highest to lowest: no chest pain (asymptomatic), atypical angina, non-anginal chest pain, typical angina} & 24\% to 14\% & log-odds -0.62, OR 0.54; range across patients -0.62 to -0.62 & 108, 42 & train; test unclear (AUC gain 0.029 [0.000, 0.065], n=61) & \textcolor{OliveGreen}{yes} & on the training data only \\
C23 & direction: maximum heart rate: 118 beats per minute to 175 beats per minute (whole model) \newline\textit{shape: accelerating; 100\% of patients same; verified} & 28\% to 21\% & log-odds -0.41, OR 0.67; range across patients -0.41 to -0.41 & 242 & train; test unclear (AUC gain 0.004 [-0.005, 0.016], n=61) & \textcolor{OliveGreen}{yes} & on the training data only \\
C24 & direction: maximum heart rate: 118 beats per minute to 175 beats per minute (direct contribution only, not through any part) \newline\textit{shape: steady; 100\% of patients same; verified} & 27\% to 21\% & log-odds -0.32, OR 0.73; range across patients -0.32 to -0.32 & 242 & train; test unclear (AUC gain 0.004 [0.000, 0.012], n=61) & \textcolor{OliveGreen}{yes} & on the training data only \\
\bottomrule
\end{tabularx}}
\par{\tiny\noindent
\begin{tabularx}{\linewidth}{@{}l>{\raggedright\arraybackslash}X>{\raggedright\arraybackslash}p{0.11\linewidth}>{\raggedright\arraybackslash}p{0.16\linewidth}>{\raggedright\arraybackslash}p{0.05\linewidth}>{\raggedright\arraybackslash}p{0.15\linewidth}l>{\raggedright\arraybackslash}p{0.08\linewidth}@{}}
\toprule
Claim & What is checked & Risk, typical patient & Estimate (log-odds; OR); range across patients & n & Test check & Verified & Words \\
\midrule
C25 & direction: age: 42 years to 66 years (whole model) \newline\textit{shape: steady; 100\% of patients same; verified} & 21\% to 25\% & log-odds 0.23, OR 1.25; range across patients 0.23 to 0.23 & 242 & train; test disagrees (AUC gain -0.002 [-0.009, 0.000], n=61) & \textcolor{OliveGreen}{yes} & slightly; on the training data only \\
C26 & magnitude: number of vessels coloured by fluoroscopy: 0 vessels to 2 vessels, typical patient \newline\textit{shape: steady (equal steps); 100\% of patients same; verified; risk saturates near 100\%} & 24\% to 98\% & log-odds 5.04, OR 154 & 242 & train; test unclear (AUC gain 0.044 [-0.038, 0.127], n=61) & \textcolor{OliveGreen}{yes} & on the training data only \\
C27 & magnitude: slope of the ST segment at peak exercise: an upsloping ST segment to a flat ST segment, typical patient \newline\textit{shape: categorical; highest to lowest: a flat ST segment, a downsloping ST segment, an upsloping ST segment} & 24\% to 82\% & log-odds 2.72, OR 15.1 & 115, 110 & train; test disagrees (AUC gain -0.013 [-0.058, 0.026], n=61) & \textcolor{OliveGreen}{yes} & on the training data only \\
C28 & magnitude: serum cholesterol: 193 mg/dL (5.0 mmol/L) to 308 mg/dL (8.0 mmol/L), typical patient \newline\textit{shape: steady; 100\% of patients same; verified} & 14\% to 44\% & log-odds 1.60, OR 4.97 & 242 & train; test disagrees (AUC gain -0.019 [-0.049, 0.002], n=61) & \textcolor{OliveGreen}{yes} & on the training data only \\
C29 & magnitude: resting blood pressure: 110 mm Hg to 155 mm Hg, typical patient \newline\textit{shape: steady; 100\% of patients same; verified} & 14\% to 41\% & log-odds 1.48, OR 4.37 & 242 & train; test unclear (AUC gain 0.002 [-0.018, 0.023], n=61) & \textcolor{OliveGreen}{yes} & on the training data only \\
C30 & magnitude: thallium stress test result: a normal thallium scan to thallium stress test result other than a normal thallium scan, a fixed defect on the thallium scan, typical patient \newline\textit{shape: categorical; highest to lowest: thallium stress test result other than a normal thallium scan, a fixed defect on the thallium scan, a normal thallium scan, a fixed defect on the thallium scan} & 24\% to 48\% & log-odds 1.09, OR 2.98 & 135, 92 & train; test unclear (AUC gain 0.005 [-0.028, 0.039], n=61) & \textcolor{OliveGreen}{yes} & on the training data only \\
C31 & structure: made of fixed\_\allowbreak{}defect\_\allowbreak{}normal\_\allowbreak{}sugar\_\allowbreak{}gate,\allowbreak{} no\_\allowbreak{}asymptomatic\_\allowbreak{}pain\_\allowbreak{}indicator,\allowbreak{} downsloping\_\allowbreak{}st\_\allowbreak{}indicator,\allowbreak{} thalach\_\allowbreak{}floor\_\allowbreak{}relay and inputs age, sex, cp, trestbps, chol, fbs, restecg, thalach, exang, oldpeak, slope, ca, thal; added &  & log-odds n/a, OR n/a & 242 & train & \textcolor{OliveGreen}{yes} &  \\
C32 & direction: fasting blood sugar: fasting blood sugar of 120 mg/dL or less to fasting blood sugar above 120 mg/dL (through this part) \newline\textit{shape: categorical; highest to lowest: fasting blood sugar above 120 mg/dL, fasting blood sugar of 120 mg/dL or less} & 24\% to 94\% (typical patient this part acts on) & log-odds 0.24, OR 1.27; range across patients 0.00 to 0.00 & 209, 33 & train; test disagrees (AUC gain -0.004 [-0.017, 0.000], n=61) & \textcolor{OliveGreen}{yes} & on the training data only \\
C33 & direction: thallium stress test result: thallium stress test result other than a fixed defect on the thallium scan to a fixed defect on the thallium scan (through this part) \newline\textit{shape: categorical; highest to lowest: thallium stress test result other than a fixed defect on the thallium scan, a fixed defect on the thallium scan} & 24\% to 1\% & log-odds -3.38, OR 0.03; range across patients -3.92 to 0.00 & 227, 15 (small n) & train; test disagrees (AUC gain -0.016 [-0.044, 0.001], n=61) & \textcolor{OliveGreen}{yes} & the model has little data on this; on the training data only \\
C34 & interaction: effect of thallium stress test result (thallium stress test result other than a fixed defect on the thallium scan to a fixed defect on the thallium scan) by fasting blood sugar: fasting blood sugar above 120 mg/dL minus fasting blood sugar of 120 mg/dL or less; stronger in fasting blood sugar of 120 mg/dL or less \newline\textit{shape: categorical; highest to lowest: thallium stress test result other than a fixed defect on the thallium scan, a fixed defect on the thallium scan} & fasting blood sugar above 120 mg/dL: 7\% to 18\%; fasting blood sugar of 120 mg/dL or less: 24\% to 2\% & log-odds difference 3.92, OR 50.3 & 33, 209, 15 (small n) & train; test agrees (AUC gain -0.004 [-0.017, 0.000], n=61); group difference on test 3.92 [3.92, 3.92] & \textcolor{OliveGreen}{yes} & the model has little data on this \\
C35 & subgroup: who this part acts on: 10 of 242 patients & 94\% to 24\% (typical patient this part acts on (risk without the part, then with it)) & share 0.04; log-odds -3.92, OR 0.02; range across patients -3.92 to -3.92 & 10, 232 (small n) & train; test disagrees (AUC gain -0.016 [-0.044, 0.000], n=61) & \textcolor{OliveGreen}{yes} & the model has little data on this; on the training data only \\
C36 & structure: its formula reproduces it on 242 training patients (largest difference 4.44089e-16) &  & log-odds 4e-16, OR n/a & 242 & train & \textcolor{OliveGreen}{yes} &  \\
\bottomrule
\end{tabularx}}
\par{\tiny\noindent
\begin{tabularx}{\linewidth}{@{}l>{\raggedright\arraybackslash}X>{\raggedright\arraybackslash}p{0.11\linewidth}>{\raggedright\arraybackslash}p{0.16\linewidth}>{\raggedright\arraybackslash}p{0.05\linewidth}>{\raggedright\arraybackslash}p{0.15\linewidth}l>{\raggedright\arraybackslash}p{0.08\linewidth}@{}}
\toprule
Claim & What is checked & Risk, typical patient & Estimate (log-odds; OR); range across patients & n & Test check & Verified & Words \\
\midrule
C37 & direction: chest pain type: no chest pain (asymptomatic) to chest pain type other than no chest pain (asymptomatic) (through this part) \newline\textit{shape: categorical; highest to lowest: no chest pain (asymptomatic), chest pain type other than no chest pain (asymptomatic)} & 24\% to 12\% & log-odds -0.85, OR 0.43; range across patients -0.85 to -0.85 & 108, 134 & train; test unclear (AUC gain 0.017 [-0.004, 0.043], n=61) & \textcolor{OliveGreen}{yes} & on the training data only \\
C38 & subgroup: who this part acts on: 134 of 242 patients & 2\% to 1\% (typical patient this part acts on (risk without the part, then with it)) & share 0.55; log-odds -0.85, OR 0.43; range across patients -0.85 to -0.85 & 134, 108 & train; test unclear (AUC gain 0.017 [-0.004, 0.043], n=61) & \textcolor{OliveGreen}{yes} & slightly; on the training data only \\
C39 & structure: its formula reproduces it on 242 training patients (largest difference 0.0) &  & log-odds 0.00, OR n/a & 242 & train & \textcolor{OliveGreen}{yes} &  \\
C40 & direction: slope of the ST segment at peak exercise: slope of the ST segment at peak exercise other than a downsloping ST segment to a downsloping ST segment (through this part) \newline\textit{shape: categorical; highest to lowest: a downsloping ST segment, slope of the ST segment at peak exercise other than a downsloping ST segment} & 24\% to 29\% & log-odds 0.30, OR 1.35; range across patients 0.30 to 0.30 & 225, 17 (small n) & train; test unclear (AUC gain 0.001 [0.000, 0.006], n=61) & \textcolor{OliveGreen}{yes} & the model has little data on this; on the training data only \\
C41 & subgroup: who this part acts on: 17 of 242 patients & 66\% to 73\% (typical patient this part acts on (risk without the part, then with it)) & share 0.07; log-odds 0.30, OR 1.35; range across patients 0.30 to 0.30 & 17, 225 (small n) & train; test unclear (AUC gain 0.001 [0.000, 0.006], n=61) & \textcolor{OliveGreen}{yes} & the model has little data on this; on the training data only \\
C42 & structure: its formula reproduces it on 242 training patients (largest difference 0.0) &  & log-odds 0.00, OR n/a & 242 & train & \textcolor{OliveGreen}{yes} &  \\
C43 & direction: maximum heart rate: 118 beats per minute to 175 beats per minute (through this part) \newline\textit{shape: hinge at 131 beats per minute; 100\% of patients same; verified; one point with C44 (median fold 131 beats per minute, typical patient 130 beats per minute, 10th-90th percentile 131 beats per minute to 131 beats per minute)} & 24\% to 23\% & log-odds -0.09, OR 0.92; range across patients -0.09 to -0.09 & 242 & train; test unclear (AUC gain 0.000 [0.000, 0.000], n=61) & \textcolor{OliveGreen}{yes} & slightly; on the training data only \\
C44 & threshold: maximum heart rate: fold at about 131 beats per minute (starts) \newline\textit{shape: hinge at 131 beats per minute; 100\% of patients same; verified; one point with C44 (median fold 131 beats per minute, typical patient 130 beats per minute, 10th-90th percentile 131 beats per minute to 131 beats per minute)} & 24\% to 23\% (typical patient with a fold) & fold 131 beats per minute (IQR 131 to 131); past it log-odds -0.10, OR 0.90; range across patients -0.10 to -0.10 & 242, 190 & train; test unclear (AUC gain 0.000 [0.000, 0.000], n=61) & \textcolor{OliveGreen}{yes} & slightly; on the training data only \\
C45 & subgroup: who this part acts on: 190 of 242 patients & 21\% to 20\% (typical patient this part acts on (risk without the part, then with it)) & share 0.79; log-odds -0.06, OR 0.95; range across patients -0.09 to -0.02 & 190, 52 & train; test unclear (AUC gain 0.000 [0.000, 0.000], n=61) & \textcolor{OliveGreen}{yes} & slightly; on the training data only \\
C46 & null: this part barely changes predictions (switched off) &  & max change 0.026 (max $|$change in predicted risk$|$) & 242 & train & \textcolor{BrickRed}{no} &  \\
C47 & structure: its formula reproduces it on 242 training patients (largest difference 0.0) &  & log-odds 0.00, OR n/a & 242 & train & \textcolor{OliveGreen}{yes} &  \\
\bottomrule
\end{tabularx}}
\paragraph{Not supported by the model.}
\begin{itemize}\setlength\itemsep{0pt}\small
\item C14: direction claim about thallium stress test result in the whole model: not supported by the model (14\% of patients have the opposite sign (allowed: 10\%)).
\item C46: null claim about maximum heart rate reached during the exercise test in part ann\_\allowbreak{}34: not supported by the model (predictions move by up to 0.026).
\end{itemize}

\subsection{Coverage}\label{app:heart-s4-c1:coverage}
Coverage is computed by the current code (the function behind the \texttt{get\_coverage} tool) under these conventions. \emph{Candidate nodes} are all nodes of the explained model state except the output nodes. A candidate node is \emph{covered} when some annotation of the hierarchy (a leaf, or a composition together with its descendants) contains it and all its enabled outgoing edges. \emph{Exact Null} parts, which have no effect on any output for any input, count as covered: nodes with no path to an output over enabled, non-zero-weight edges (disconnected inputs and dead ends), and edges that have zero weight or touch such a node. Compositional coverage is the fraction of required compositions present: every annotation except a single global root containing every output must have a parent; a parentless, childless annotation made only of Null nodes is exempt.

\begin{itemize}
\item \textbf{Null, disconnected input (6):} \texttt{-12} (restecg:0), \texttt{-14} (restecg:2), \texttt{-16} (exang:0), \texttt{-24} (thal:4), \texttt{-26} (thal:7), \texttt{-5} (cp:2).
\item \textbf{Uncovered nodes:} none.
\item \textbf{Uncovered edges:} none.
\item \textbf{Missing compositions:} none.
\end{itemize}

\subsection{Annotation hierarchy}\label{app:heart-s4-c1:hierarchy}
The hierarchy has 1 root; the global root is heart\_\allowbreak{}disease\_\allowbreak{}logit. Tags (A1, A2, \ldots) are used in the figures.
\begingroup\parindent=0pt\parskip=1pt
\hspace*{0.0em}$\triangleright$~A1: heart\_\allowbreak{}disease\_\allowbreak{}logit (composition, \cref{app:heart-s4-c1:ann-heart-disease-logit})\par
\hspace*{1.5em}$\triangleright$~A2: fixed\_\allowbreak{}defect\_\allowbreak{}normal\_\allowbreak{}sugar\_\allowbreak{}gate (leaf, \cref{app:heart-s4-c1:ann-fixed-defect-normal-sugar-gate})\par
\hspace*{1.5em}$\triangleright$~A3: no\_\allowbreak{}asymptomatic\_\allowbreak{}pain\_\allowbreak{}indicator (leaf, \cref{app:heart-s4-c1:ann-no-asymptomatic-pain-indicator})\par
\hspace*{1.5em}$\triangleright$~A4: downsloping\_\allowbreak{}st\_\allowbreak{}indicator (leaf, \cref{app:heart-s4-c1:ann-downsloping-st-indicator})\par
\hspace*{1.5em}$\triangleright$~A5: thalach\_\allowbreak{}floor\_\allowbreak{}relay (leaf, \cref{app:heart-s4-c1:ann-thalach-floor-relay})\par
\endgroup

\subsection{Network figures}\label{app:heart-s4-c1:figures}
\begin{figure}[htbp]
\centering
\input{\explaneatcertdir{heart-s4-c1}heart-s4-c1-fig-full.tex}
\caption{heart-s4-c1: full view of the explained model state.}
\label{app:heart-s4-c1:fig-full}
\end{figure}
\begin{figure}[htbp]
\centering
\input{\explaneatcertdir{heart-s4-c1}heart-s4-c1-fig-level-0.tex}
\caption{heart-s4-c1: collapsed view at hierarchy level 0 (fully recomposed).}
\label{app:heart-s4-c1:fig-level-0}
\end{figure}
\begin{figure}[htbp]
\centering
\input{\explaneatcertdir{heart-s4-c1}heart-s4-c1-fig-level-1.tex}
\caption{heart-s4-c1: collapsed view at hierarchy level 1.}
\label{app:heart-s4-c1:fig-level-1}
\end{figure}

\subsection{Refactoring operations}\label{app:heart-s4-c1:operations}
The explained model state is the phenotype after these operations, replayed in order. All are function-preserving; the largest change in the model output over the 242 training rows is $3.33\times 10^{-16}$ (tolerance $1.00\times 10^{-9}$).
\begin{itemize}
\item \#26: split \texttt{-10} into \texttt{-10\_\allowbreak{}a}, \texttt{-10\_\allowbreak{}b}.
\item \#27: split \texttt{-11} into \texttt{-11\_\allowbreak{}a}, \texttt{-11\_\allowbreak{}b}.
\item \#28: split \texttt{-15} into \texttt{-15\_\allowbreak{}a}, \texttt{-15\_\allowbreak{}b}.
\item \#29: split \texttt{-21} into \texttt{-21\_\allowbreak{}a}, \texttt{-21\_\allowbreak{}b}.
\item \#30: split \texttt{-7} into \texttt{-7\_\allowbreak{}a}, \texttt{-7\_\allowbreak{}b}.
\end{itemize}

\subsection{Annotations}\label{app:heart-s4-c1:annotations}
\subsubsection{\texorpdfstring{A1: heart\_\allowbreak{}disease\_\allowbreak{}logit}{A1: heart disease logit}}\label{app:heart-s4-c1:ann-heart-disease-logit}
% annotation id ann_35, operation seq 35
\textbf{Kind:} Composition annotation (depth 0); children A2 (\cref{app:heart-s4-c1:ann-fixed-defect-normal-sugar-gate}), A3 (\cref{app:heart-s4-c1:ann-no-asymptomatic-pain-indicator}), A4 (\cref{app:heart-s4-c1:ann-downsloping-st-indicator}), A5 (\cref{app:heart-s4-c1:ann-thalach-floor-relay}).

\begin{itemize}
\item \textbf{Entry nodes:} \texttt{-1} (age), \texttt{-2} (sex:0), \texttt{-3} (sex:1), \texttt{-4} (cp:1), \texttt{-6} (cp:3), \texttt{-8} (trestbps), \texttt{-9} (chol), \texttt{-13} (restecg:1), \texttt{-17} (exang:1), \texttt{-18} (oldpeak), \texttt{-19} (slope:1), \texttt{-20} (slope:2), \texttt{-22} (ca), \texttt{-23} (thal:3), \texttt{-10\_\allowbreak{}a} (fbs:0\_\allowbreak{}a), \texttt{-11\_\allowbreak{}a} (fbs:1\_\allowbreak{}a), \texttt{-15\_\allowbreak{}a} (thalach\_\allowbreak{}a), \texttt{-21\_\allowbreak{}a} (slope:3\_\allowbreak{}a), \texttt{-7\_\allowbreak{}a} (cp:4\_\allowbreak{}a), \texttt{-26} (thal:7), \texttt{-24} (thal:4), \texttt{-16} (exang:0), \texttt{-14} (restecg:2), \texttt{-12} (restecg:0), \texttt{-5} (cp:2), \texttt{-10\_\allowbreak{}b} (fbs:0\_\allowbreak{}b), \texttt{-11\_\allowbreak{}b} (fbs:1\_\allowbreak{}b), \texttt{-25} (thal:6), \texttt{-7\_\allowbreak{}b} (cp:4\_\allowbreak{}b), \texttt{-21\_\allowbreak{}b} (slope:3\_\allowbreak{}b), \texttt{-15\_\allowbreak{}b} (thalach\_\allowbreak{}b).
\item \textbf{Exit nodes:} \texttt{0}.
\item \textbf{Internal nodes:} none.
\item \textbf{Edges (weight):} \texttt{-1}$\to$\texttt{0}~(0.0855), \texttt{-2}$\to$\texttt{0}~(-0.559), \texttt{-3}$\to$\texttt{0}~(0.899), \texttt{-4}$\to$\texttt{0}~(-0.679), \texttt{-6}$\to$\texttt{0}~(-0.749), \texttt{-8}$\to$\texttt{0}~(0.596), \texttt{-9}$\to$\texttt{0}~(0.671), \texttt{-13}$\to$\texttt{0}~(0.0618), \texttt{-17}$\to$\texttt{0}~(0.282), \texttt{-18}$\to$\texttt{0}~(0.491), \texttt{-19}$\to$\texttt{0}~(-0.238), \texttt{-20}$\to$\texttt{0}~(1.11), \texttt{-22}$\to$\texttt{0}~(2.25), \texttt{-23}$\to$\texttt{0}~(-0.541), \texttt{-10\_\allowbreak{}a}$\to$\texttt{0}~(-0.0936), \texttt{-11\_\allowbreak{}a}$\to$\texttt{0}~(-0.583), \texttt{-15\_\allowbreak{}a}$\to$\texttt{0}~(-0.125), \texttt{-21\_\allowbreak{}a}$\to$\texttt{0}~(0.0688), \texttt{-7\_\allowbreak{}a}$\to$\texttt{0}~(0.308), \texttt{2065}$\to$\texttt{0}~(-1.08), \texttt{1754}$\to$\texttt{0}~(-0.899), \texttt{3504}$\to$\texttt{0}~(0.114), \texttt{707}$\to$\texttt{0}~(-0.0469).
\item \textbf{With its descendants:} 36 nodes.
\item \textbf{Covers} (coverage, own and as part of the hierarchy): 35 candidate nodes.
\end{itemize}

\textbf{Hypothesis.} Root: p = sigmoid(z), z = 0.466 + LINEAR CORE - 1.077*\allowbreak{}fixed\_\allowbreak{}defect\_\allowbreak{}normal\_\allowbreak{}sugar\_\allowbreak{}gate - 0.899*\allowbreak{}no\_\allowbreak{}asymptomatic\_\allowbreak{}pain\_\allowbreak{}indicator + 0.114*\allowbreak{}downsloping\_\allowbreak{}st\_\allowbreak{}indicator - 0.047*\allowbreak{}thalach\_\allowbreak{}floor\_\allowbreak{}relay. LINEAR CORE (19 direct input-$>$output edges, z-scored inputs): ca +2.246, slope:2 +1.115, sex:1 +0.899, cp:3 -0.749, cp:1 -0.679, chol +0.671, trestbps +0.596, fbs:1 -0.583, sex:0 -0.559, thal:3 -0.541, oldpeak +0.491, cp:4 +0.308, exang:1 +0.282, slope:1 -0.238, thalach -0.125, fbs:0 -0.094, age +0.085, slope:3 +0.069, restecg:1 +0.062. Six inputs are unconnected (cp:2, restecg:0, restecg:2, exang:0, thal:4, thal:7; zero-effect entries of this root). thal:6 has no direct edge; it acts only through the fixed-defect gate. COMPOSITION: purely additive in the logit. Children: (1) fixed\_\allowbreak{}defect\_\allowbreak{}normal\_\allowbreak{}sugar\_\allowbreak{}gate,\allowbreak{} the one genuine interaction: an AND of fixed thallium defect and normal fasting sugar that lowers the logit by 3.92 for 10 patients; (2) no\_\allowbreak{}asymptomatic\_\allowbreak{}pain\_\allowbreak{}indicator and (3) downsloping\_\allowbreak{}st\_\allowbreak{}indicator,\allowbreak{} binary ReLUs exactly affine on their one-hot inputs (re-parameterisations of the cp:4 and slope:3 linear terms; cp:4 total +1.47 logit, slope:3 total +0.57); (4) thalach\_\allowbreak{}floor\_\allowbreak{}relay, a near-inert hinge on max heart rate. The children share no hidden nodes; they are coupled only through inputs that also have direct edges (fbs:0, fbs:1, cp:4, slope:3, thalach). EVIDENCE (train, n=242): full model AUC 0.9423, acc 0.888 (compute\_\allowbreak{}performance)\allowbreak{}. Linear core alone (direct edges only) AUC 0.9328; logit SD 4.55 (core) vs 0.86 (all hidden readouts). Zero-ablating one child at a time: fixed-defect gate 0.9343, cp:4 indicator 0.9411, thalach relay 0.9416, slope:3 indicator 0.9421. Replacing all hidden readouts with one least-squares linear fit in their six inputs gives 0.9371; an L2 logistic regression on all 26 inputs gives 0.9355. Mean $|$SHAP$|$ (train, all 242, probability scale): ca 0.172, slope:2 0.099, sex:1 0.069, cp:4 0.068, cp:3 0.056, thal:3 0.049, sex:0 0.042, chol 0.042. Interpretation: essentially logistic regression with one meaningful nonlinearity, the fixed-defect/\allowbreak{}normal-sugar gate, which cancels the core's over-prediction for that small group (10 patients). Counter-directional signs reported, not explained: fbs:1 -0.58 (high fasting sugar lowers the logit) and the gate itself making a fixed defect protective for normal-sugar patients. Same mechanism family as the 10-seed heart-s4 record (key 9650, gen 38): the same gate and linear core (ca +2.16 there), but here a cp:4 indicator and a slope:3 indicator replace that genome's cp:2 indicator, and the thalach relay feeds the output directly.

\textbf{Closed form} (children as named functions):
{\small\begingroup\allowdisplaybreaks
\begin{align*}
y_{0} &= \sigma\bigl(z_{0}\bigr) \\
z_{0} &= 0.466 + 0.0688 x_{slope:3\_a} + 0.282 x_{exang:1} + 0.308 x_{cp:4\_a} \\
&\qquad + 0.899 x_{sex:1} + 0.491 x_{oldpeak} \\
&\qquad + 0.114 \operatorname{downsloping\_st\_indicator}\bigl(x_{slope:3\_b}\bigr) + 1.11 x_{slope:2} \\
&\qquad + 2.25 x_{ca} + 0.596 x_{trestbps} + 0.671 x_{chol} + 0.0855 x_{age} + 0.0618 x_{restecg:1} \\
&\qquad - 0.899 \operatorname{no\_asymptomatic\_pain\_indicator}\bigl(x_{cp:4\_b}\bigr) \\
&\qquad - 0.542 x_{thal:3} - 0.583 x_{fbs:1\_a} - 0.0936 x_{fbs:0\_a} - 0.749 x_{cp:3} \\
&\qquad - 0.0469 \operatorname{thalach\_floor\_relay}\bigl(x_{thalach\_b}\bigr) - 0.238 x_{slope:1} \\
&\qquad - 1.08 \operatorname{fixed\_defect\_normal\_sugar\_gate}\bigl(x_{fbs:0\_b}, x_{fbs:1\_b}, \\
&\qquad x_{thal:6}\bigr) - 0.559 x_{sex:0} - 0.679 x_{cp:1} - 0.125 x_{thalach\_a}
\end{align*}\endgroup}

Expanded over the children:
{\small\begingroup\allowdisplaybreaks
\begin{align*}
y_{0} &= \sigma\bigl(z_{0}\bigr) \\
z_{0} &= 0.466 + 0.0688 x_{slope:3\_a} + 0.282 x_{exang:1} + 0.308 x_{cp:4\_a} \\
&\qquad + 0.899 x_{sex:1} + 0.491 x_{oldpeak} + 0.114 \operatorname{relu}\bigl(-1.02 \\
&\qquad + 1.01 x_{slope:3\_b}\bigr) + 1.11 x_{slope:2} + 2.25 x_{ca} + 0.596 x_{trestbps} \\
&\qquad + 0.671 x_{chol} + 0.0855 x_{age} + 0.0618 x_{restecg:1} \\
&\qquad - 0.899 \operatorname{relu}\bigl(-0.0387 - 1.1 x_{cp:4\_b}\bigr) - 0.542 x_{thal:3} \\
&\qquad - 0.583 x_{fbs:1\_a} - 0.0936 x_{fbs:0\_a} - 0.749 x_{cp:3} \\
&\qquad - 0.0469 \operatorname{relu}\bigl(0.812 + 0.959 x_{thalach\_b}\bigr) - 0.238 x_{slope:1} \\
&\qquad - 1.08 \operatorname{relu}\bigl(-0.703 + 1.58 x_{fbs:0\_b} + 0.885 x_{thal:6} \\
&\qquad - 0.683 x_{fbs:1\_b}\bigr) - 0.559 x_{sex:0} - 0.679 x_{cp:1} - 0.125 x_{thalach\_a}
\end{align*}\endgroup}

where $y_{0}$ is node \texttt{0}; $z_{i}$, where shown, is the pre-activation of a sigmoid output $y_{i}$; weights and biases are shown to 3 significant figures.
 On the training split (n = 242), the displayed closed form differs from the forward pass by at most 0.00419, and the exact function by at most $2.78\times 10^{-16}$.

\textbf{Evidence} (2 items; full payloads in the JSON record):
\begin{itemize}
\item \textsc{analytical}, \texttt{analytical\_\allowbreak{}formula}; data split: n/a (analytical). Root closed form (get\_\allowbreak{}formula force=true); output is sigmoid.
\item \textsc{empirical}, \texttt{empirical\_\allowbreak{}viz}; data split: train. Train-split performance, ablations, linear replacements and SHAP.
\end{itemize}

\textbf{Authorship:} drafted by AI-assisted draft on 2026-11-10; verified by nobody (unverified).

\subsubsection{\texorpdfstring{A2: fixed\_\allowbreak{}defect\_\allowbreak{}normal\_\allowbreak{}sugar\_\allowbreak{}gate}{A2: fixed defect normal sugar gate}}\label{app:heart-s4-c1:ann-fixed-defect-normal-sugar-gate}
% annotation id ann_31, operation seq 31
\textbf{Kind:} Leaf annotation (depth 1); parent A1 (\cref{app:heart-s4-c1:ann-heart-disease-logit}).

\begin{itemize}
\item \textbf{Entry nodes:} \texttt{-10\_\allowbreak{}b} (fbs:0\_\allowbreak{}b), \texttt{-11\_\allowbreak{}b} (fbs:1\_\allowbreak{}b), \texttt{-25} (thal:6).
\item \textbf{Exit nodes:} \texttt{2065}.
\item \textbf{Internal nodes:} none.
\item \textbf{Edges (weight):} \texttt{-10\_\allowbreak{}b}$\to$\texttt{2065}~(1.58), \texttt{-11\_\allowbreak{}b}$\to$\texttt{2065}~(-0.683), \texttt{-25}$\to$\texttt{2065}~(0.885).
\item \textbf{Covers} (coverage, own and as part of the hierarchy): 3 candidate nodes.
\end{itemize}

\textbf{Hypothesis.} Fixed-defect AND normal-sugar gate (lowers risk). 2065 = ReLU(-0.703 + 1.576*fbs:0 - 0.683*fbs:1 + 0.885*thal:6), read out -1.077. Hypothesis: a discrete AND detector that is on only for patients with a fixed thallium defect (thal 6) and fasting blood sugar $<$= 120 mg/dL, and then lowers their logit by 3.92, cancelling the linear core's over-calling of this group (thal:6 has no direct edge; the core gives these patients a mean predicted probability of 0.83 while their disease rate is 50\%). EVIDENCE (train, n=242): exact on-data reduction 2065 = 3.639*[thal:6=1 AND fbs=0] (fbs=1 drives the pre-activation to $<$= -6.39). Active for exactly the 10 patients in that group; disease rate 5/10. Zero-ablating it lowers train AUC 0.9423 -$>$ 0.9343, the largest nonlinear effect; a linear refit in its three inputs recovers only to 0.937 (R2 0.69), because an AND cannot be written additively. Mean $|$sensitivity$|$ fbs:0 0.065 $>$ thal:6 0.037 $>$ fbs:1 0.028 (small because the gate is off for 96\%). Supported, but the correction rests on 10 training patients and may be partly noise-fitted. Same gate as the 10-seed heart-s4 record (gen 38), now also fed by fbs:0.

\textbf{Closed form:}
{\small\begingroup\allowdisplaybreaks
\begin{align*}
y_{0} &= \operatorname{relu}\bigl(-0.703 + 1.58 x_{fbs:0\_b} + 0.885 x_{thal:6} \\
&\qquad - 0.683 x_{fbs:1\_b}\bigr)
\end{align*}\endgroup}

where $y_{0}$ is node \texttt{2065}; $z_{i}$, where shown, is the pre-activation of a sigmoid output $y_{i}$; weights and biases are shown to 3 significant figures.
 On the training split (n = 242), the displayed closed form differs from the forward pass by at most $4.34\times 10^{-4}$, and the exact function by at most 0.

\textbf{Evidence} (2 items; full payloads in the JSON record):
\begin{itemize}
\item \textsc{analytical}, \texttt{analytical\_\allowbreak{}formula}; data split: n/a (analytical). Closed form from get\_\allowbreak{}formula with the on-data simplification.
\item \textsc{empirical}, \texttt{empirical\_\allowbreak{}viz} (sensitivity (max\_\allowbreak{}samples 100000)); data split: train. Train-split sensitivity, activity, outcome rates and ablation.
\end{itemize}

\textbf{Authorship:} drafted by AI-assisted draft on 2026-11-10; verified by nobody (unverified).

\subsubsection{\texorpdfstring{A3: no\_\allowbreak{}asymptomatic\_\allowbreak{}pain\_\allowbreak{}indicator}{A3: no asymptomatic pain indicator}}\label{app:heart-s4-c1:ann-no-asymptomatic-pain-indicator}
% annotation id ann_32, operation seq 32
\textbf{Kind:} Leaf annotation (depth 1); parent A1 (\cref{app:heart-s4-c1:ann-heart-disease-logit}).

\begin{itemize}
\item \textbf{Entry nodes:} \texttt{-7\_\allowbreak{}b} (cp:4\_\allowbreak{}b).
\item \textbf{Exit nodes:} \texttt{1754}.
\item \textbf{Internal nodes:} none.
\item \textbf{Edges (weight):} \texttt{-7\_\allowbreak{}b}$\to$\texttt{1754}~(-1.1).
\item \textbf{Covers} (coverage, own and as part of the hierarchy): 1 candidate node.
\end{itemize}

\textbf{Hypothesis.} Not-asymptomatic chest-pain indicator (lowers risk). 1754 = ReLU(-0.039 - 1.098*cp:4), read out -0.899. Hypothesis: a 0/1 indicator for chest pain other than asymptomatic (typical, atypical or non-anginal pain); because cp:4 is binary it is exactly affine on the data, so it only adds to the cp:4 direct edge (+0.308 per z-unit): together, asymptomatic chest pain raises the logit by 1.47 relative to the other pain types. EVIDENCE (train, n=242): on for exactly the 134 patients without asymptomatic pain (value 0.947, logit -0.852); disease rate 22\% when on vs 71\% when off. A linear replacement in cp:4 reproduces it exactly (R2 1.00, AUC unchanged 0.9423); zero-ablation (which also removes cp:4's main effect) gives 0.9411. Supported: a re-parameterisation of a linear cp:4 term, no interaction.

\textbf{Closed form:}
{\small\begingroup\allowdisplaybreaks
\begin{align*}
y_{0} &= \operatorname{relu}\bigl(-0.0387 - 1.1 x_{cp:4\_b}\bigr)
\end{align*}\endgroup}

where $y_{0}$ is node \texttt{1754}; $z_{i}$, where shown, is the pre-activation of a sigmoid output $y_{i}$; weights and biases are shown to 3 significant figures.
 On the training split (n = 242), the displayed closed form differs from the forward pass by at most 0.00173, and the exact function by at most 0.

\textbf{Evidence} (2 items; full payloads in the JSON record):
\begin{itemize}
\item \textsc{analytical}, \texttt{analytical\_\allowbreak{}formula}; data split: n/a (analytical). Closed form from get\_\allowbreak{}formula with the on-data simplification.
\item \textsc{empirical}, \texttt{empirical\_\allowbreak{}viz} (sensitivity (max\_\allowbreak{}samples 100000)); data split: train. Train-split sensitivity, activity, outcome rates and ablation.
\end{itemize}

\textbf{Authorship:} drafted by AI-assisted draft on 2026-11-10; verified by nobody (unverified).

\subsubsection{\texorpdfstring{A4: downsloping\_\allowbreak{}st\_\allowbreak{}indicator}{A4: downsloping st indicator}}\label{app:heart-s4-c1:ann-downsloping-st-indicator}
% annotation id ann_33, operation seq 33
\textbf{Kind:} Leaf annotation (depth 1); parent A1 (\cref{app:heart-s4-c1:ann-heart-disease-logit}).

\begin{itemize}
\item \textbf{Entry nodes:} \texttt{-21\_\allowbreak{}b} (slope:3\_\allowbreak{}b).
\item \textbf{Exit nodes:} \texttt{3504}.
\item \textbf{Internal nodes:} none.
\item \textbf{Edges (weight):} \texttt{-21\_\allowbreak{}b}$\to$\texttt{3504}~(1.01).
\item \textbf{Covers} (coverage, own and as part of the hierarchy): 1 candidate node.
\end{itemize}

\textbf{Hypothesis.} Downsloping-ST indicator (raises risk slightly). 3504 = ReLU(-1.020 + 1.010*slope:3), read out +0.114. Hypothesis: a 0/1 indicator for a downsloping peak-exercise ST segment; exactly affine on the binary input, adding +0.30 logit to slope:3's direct edge (total +0.57 logit for downsloping vs not). EVIDENCE (train, n=242): on for exactly the 17 downsloping patients; disease rate 47\% vs 44\%. Linear replacement exact (R2 1.00); zero-ablation 0.9423 -$>$ 0.9421. Supported but nearly inert; rests on 17 patients.

\textbf{Closed form:}
{\small\begingroup\allowdisplaybreaks
\begin{align*}
y_{0} &= \operatorname{relu}\bigl(-1.02 + 1.01 x_{slope:3\_b}\bigr)
\end{align*}\endgroup}

where $y_{0}$ is node \texttt{3504}; $z_{i}$, where shown, is the pre-activation of a sigmoid output $y_{i}$; weights and biases are shown to 3 significant figures.
 On the training split (n = 242), the displayed closed form differs from the forward pass by at most $1.95\times 10^{-4}$, and the exact function by at most 0.

\textbf{Evidence} (2 items; full payloads in the JSON record):
\begin{itemize}
\item \textsc{analytical}, \texttt{analytical\_\allowbreak{}formula}; data split: n/a (analytical). Closed form from get\_\allowbreak{}formula with the on-data simplification.
\item \textsc{empirical}, \texttt{empirical\_\allowbreak{}viz} (sensitivity (max\_\allowbreak{}samples 100000)); data split: train. Train-split sensitivity, activity, outcome rates and ablation.
\end{itemize}

\textbf{Authorship:} drafted by AI-assisted draft on 2026-11-10; verified by nobody (unverified).

\subsubsection{\texorpdfstring{A5: thalach\_\allowbreak{}floor\_\allowbreak{}relay}{A5: thalach floor relay}}\label{app:heart-s4-c1:ann-thalach-floor-relay}
% annotation id ann_34, operation seq 34
\textbf{Kind:} Leaf annotation (depth 1); parent A1 (\cref{app:heart-s4-c1:ann-heart-disease-logit}).

\begin{itemize}
\item \textbf{Entry nodes:} \texttt{-15\_\allowbreak{}b} (thalach\_\allowbreak{}b).
\item \textbf{Exit nodes:} \texttt{707}.
\item \textbf{Internal nodes:} none.
\item \textbf{Edges (weight):} \texttt{-15\_\allowbreak{}b}$\to$\texttt{707}~(0.959).
\item \textbf{Covers} (coverage, own and as part of the hierarchy): 1 candidate node.
\end{itemize}

\textbf{Hypothesis.} Max-heart-rate floor relay (lowers risk a little). 707 = ReLU(0.812 + 0.959*thalach), read out -0.047. Hypothesis: a hinge that adds a small extra protective slope in maximum heart rate achieved for all but the lowest \textasciitilde{}21\% of thalach values (kink at z-thalach -0.85); combined with the direct edge (-0.125 per z) the thalach effect is -0.125 per z below the kink and -0.170 per z above it. EVIDENCE (train, n=242): active for 78.5\%; disease rate 37\% when active vs 69\% when off; readout logit SD only 0.035. Zero-ablating it lowers train AUC 0.9423 -$>$ 0.9416; a straight-line refit gives 0.9420 (R2 0.90). Supported: a near-linear, nearly inert extension of the thalach term.

\textbf{Closed form:}
{\small\begingroup\allowdisplaybreaks
\begin{align*}
y_{0} &= \operatorname{relu}\bigl(0.812 + 0.959 x_{thalach\_b}\bigr)
\end{align*}\endgroup}

where $y_{0}$ is node \texttt{707}; $z_{i}$, where shown, is the pre-activation of a sigmoid output $y_{i}$; weights and biases are shown to 3 significant figures.
 On the training split (n = 242), the displayed closed form differs from the forward pass by at most $8.60\times 10^{-4}$, and the exact function by at most 0.

\textbf{Evidence} (2 items; full payloads in the JSON record):
\begin{itemize}
\item \textsc{analytical}, \texttt{analytical\_\allowbreak{}formula}; data split: n/a (analytical). Closed form from get\_\allowbreak{}formula with the on-data simplification.
\item \textsc{empirical}, \texttt{empirical\_\allowbreak{}viz} (sensitivity (max\_\allowbreak{}samples 100000)); data split: train. Train-split sensitivity, activity, outcome rates and ablation.
\end{itemize}

\textbf{Authorship:} drafted by AI-assisted draft on 2026-11-10; verified by nobody (unverified).

\subsection{Audit against the definitions}\label{app:heart-s4-c1:audit}
Automatic findings on the stored explanation (codes as in the JSON record's \texttt{issues}); they flag what an examiner should check, not verdicts.
\begin{itemize}
\item \textsc{warning} \texttt{unverified}: No human verification is recorded for this explanation.
\end{itemize}
\end{drafted}
